\chapter{Tight Frames}
\label{ch:tight-frames}
\fancyhead[R]{\small\textit{Chapter 4: Tight Frames}}
\section{Overview: Closing the Loop on Tightness}
\label{sec:ch4-overview}

Three separate promises from Chapters~1--3 converge in this chapter.
Chapter~1 conjectured, from a single picture, that the vertices of an
equilateral triangle form a tight frame for $\R^2$. Chapter~2 verified
this by direct computation for $m=3$ and asked you to guess, in Lab
Exercise~1.4, whether the same holds for a square. Chapter~3 confirmed
the square case too ($A=B=2$) and asked you, in Lab Exercise~3.6, to
test pentagons, hexagons, heptagons, and octagons, and to conjecture a
general pattern. By now the pattern should be unmistakable: \emph{every}
regular polygon, of any number of sides $m\ge 2$, gives a tight frame for
$\R^2$. This chapter proves that conjecture in full generality
(Theorem~\ref{thm:regular-polygon-tight}), and along the way develops the
algebraic mechanism---cancellation of trigonometric cross-terms via roots
of unity---that explains \emph{why} it is true, not merely \emph{that}
it is true.

We also deliver on a second promise: Chapter~2 noted that Parseval
frames need not be orthonormal bases, and deferred examples to this
chapter. Section~\ref{sec:parseval-from-projection} gives a general,
constructive method---projecting an orthonormal basis of a larger space
down into a subspace---that produces Parseval frames with as much
redundancy as you like, and explains geometrically why this construction
always works.

\section{Tight Frames: Definitions Revisited}
\label{sec:tight-revisited}

Recall Definition~2.4.1: a frame $\{f_i\}_{i=1}^m$ for $V$ is \emph{tight}
if its frame bounds satisfy $A = B$, and \emph{Parseval} if additionally
$A=1$. Chapter~3 gave the matrix characterization that organizes this
entire chapter:

\begin{quote}
\itshape A frame is tight with bound $A$ if and only if its frame
operator satisfies $S = AI$ (Theorem~3.5.1(iii)).
\end{quote}

This single equation, $S = AI$, is the lens through which we will view
every result in this chapter. Unfolding $S = \sum_i f_if_i^* = AI$ gives
the following equivalent, and extremely useful, restatement.

\begin{proposition}[Tightness as a resolution of the identity]
\label{prop:tightness-resolution}
A frame $\{f_i\}_{i=1}^m$ for $V=\F^n$ is tight with bound $A$ if and
only if
\begin{equation}
\label{eq:resolution-of-identity}
  \sum_{i=1}^m f_i f_i^* = A I_n.
\end{equation}
Equivalently, for every pair of standard basis vectors $e_j, e_k$,
\begin{equation}
\label{eq:entrywise-tight}
  \sum_{i=1}^m (f_i)_j \overline{(f_i)_k} =
  \begin{cases} A, & j = k, \\ 0, & j \ne k. \end{cases}
\end{equation}
\end{proposition}

\begin{proof}
This is immediate from $S = \sum_i f_if_i^* = AI$ together with the fact
that the $(j,k)$ entry of $f_if_i^*$ is $(f_i)_j\overline{(f_i)_k}$, and
the $(j,k)$ entry of $AI$ is $A$ when $j=k$ and $0$ otherwise.
\end{proof}

\begin{remark*}
Equation~\eqref{eq:resolution-of-identity} is sometimes called a
\emph{resolution of the identity}: the frame vectors' outer products sum
exactly to (a multiple of) the identity. This is the precise algebraic
form of the informal phrase from Chapter~1, ``every direction is treated
equally.'' Equation~\eqref{eq:entrywise-tight} is the same statement
written entrywise, and it is this entrywise form that we will verify
directly, by trigonometric computation, in the regular polygon theorem
below: the off-diagonal cancellation in
Example~2.2.5---the vanishing of the cross term
$-\tfrac{\sqrt3}{2}x_1x_2+\tfrac{\sqrt3}{2}x_1x_2$---is exactly an
instance of the $j\ne k$ case of~\eqref{eq:entrywise-tight}.
\end{remark*}

\section{The Regular Polygon Theorem}
\label{sec:regular-polygon}

We now state and prove the general result toward which Chapters~1--3
have been building.

\begin{theorem}[Regular polygons are tight frames]
\label{thm:regular-polygon-tight}
Let $m \ge 2$ be an integer, and let
\[
  f_k = \bigl(\cos(2\pi k/m),\, \sin(2\pi k/m)\bigr), \qquad k = 0, 1, \ldots, m-1,
\]
be the $m$ vectors pointing toward the vertices of a regular $m$-gon
inscribed in the unit circle in $\R^2$. Then $\{f_k\}_{k=0}^{m-1}$ is a
tight frame for $\R^2$ with frame bound $A = m/2$.
\end{theorem}

We give two proofs. The first is a direct trigonometric computation,
generalizing exactly the calculation of Example~2.2.5. The second is
shorter and more conceptual, using complex exponentials and roots of
unity; it is the ``systematic'' technique promised in Chapter~2; and it
will generalize far more easily to other constructions later in the
course.

\begin{proof}[Proof 1: Direct trigonometric computation]
Let $\theta_k = 2\pi k/m$, so $f_k = (\cos\theta_k, \sin\theta_k)$. By
Proposition~\ref{prop:tightness-resolution}, it suffices to verify
\eqref{eq:entrywise-tight} for $A = m/2$, i.e.\ to show
\[
  \sum_{k=0}^{m-1}\cos^2\theta_k = \frac{m}{2}, \qquad
  \sum_{k=0}^{m-1}\sin^2\theta_k = \frac{m}{2}, \qquad
  \sum_{k=0}^{m-1}\cos\theta_k\sin\theta_k = 0.
\]
Using the double-angle identities $\cos^2\theta = \tfrac12(1+\cos2\theta)$,
$\sin^2\theta = \tfrac12(1-\cos2\theta)$, and
$\cos\theta\sin\theta = \tfrac12\sin2\theta$:
\[
  \sum_{k=0}^{m-1}\cos^2\theta_k = \frac{m}{2} + \frac12\sum_{k=0}^{m-1}\cos(2\theta_k),
  \qquad
  \sum_{k=0}^{m-1}\sin^2\theta_k = \frac{m}{2} - \frac12\sum_{k=0}^{m-1}\cos(2\theta_k),
\]
\[
  \sum_{k=0}^{m-1}\cos\theta_k\sin\theta_k = \frac12\sum_{k=0}^{m-1}\sin(2\theta_k).
\]
So it suffices to show that both sums $\sum_{k=0}^{m-1}\cos(2\theta_k)$
and $\sum_{k=0}^{m-1}\sin(2\theta_k)$ vanish, where $2\theta_k = 4\pi
k/m$. These are the real and imaginary parts of the geometric series
\[
  \sum_{k=0}^{m-1} e^{i \cdot 4\pi k/m} = \sum_{k=0}^{m-1} \bigl(e^{i4\pi/m}\bigr)^k.
\]
If $m \ne 2$, then $\zeta := e^{i4\pi/m} \ne 1$ (since $4\pi/m$ is not a
multiple of $2\pi$ for $m \ge 3$), and the finite geometric series formula
gives
\[
  \sum_{k=0}^{m-1}\zeta^k = \frac{\zeta^m - 1}{\zeta - 1} = \frac{e^{i4\pi}-1}{\zeta-1} = \frac{1-1}{\zeta-1} = 0,
\]
using $e^{i4\pi}=1$. Taking real and imaginary parts gives
$\sum_k\cos(2\theta_k)=0$ and $\sum_k\sin(2\theta_k)=0$, as required.
(The case $m=2$ gives $f_0=(1,0)$, $f_1=(-1,0)$, which one checks
directly: $\sum\cos^2\theta_k = 1+1=2=m/2\cdot 2$... we revisit $m=2$
explicitly in Remark~\ref{rmk:m-equals-2} below, since the geometric
series argument needs $\zeta\ne1$.) This proves
$\sum_k\cos^2\theta_k=\sum_k\sin^2\theta_k=m/2$ and
$\sum_k\cos\theta_k\sin\theta_k=0$, which is exactly
\eqref{eq:entrywise-tight} with $A=m/2$.
\end{proof}

\begin{proof}[Proof 2: Roots of unity (the systematic technique)]
Identify $\R^2$ with $\C$ via $(\cos\theta,\sin\theta)\leftrightarrow
e^{i\theta}$. Let $\omega = e^{2\pi i/m}$, a primitive $m$-th root of
unity, so that $f_k \leftrightarrow \omega^k$ for $k=0,\ldots,m-1$. We
argue directly with the frame operator acting on a vector $w \in \C$
(thought of as a vector in $\R^2$ under the identification above). Since
$f_k\leftrightarrow\omega^k$, the real inner product
$\inner{w}{f_k}$ corresponds to $\Real(\bar w\,\omega^k)$, so
\[
  \sum_{k=0}^{m-1}|\inner{w}{f_k}|^2
  = \sum_{k=0}^{m-1}\Real(\bar w\,\omega^k)^2
  = \sum_{k=0}^{m-1}\frac{(\bar w\,\omega^k+w\,\bar\omega^k)^2}{4}.
\]
Expanding the square,
\[
  (\bar w\,\omega^k+w\,\bar\omega^k)^2
  = \bar w^2\omega^{2k} + 2|w|^2 + w^2\bar\omega^{2k},
\]
so
\[
  \sum_{k=0}^{m-1}|\inner{w}{f_k}|^2
  = \frac{\bar w^2}{4}\sum_{k=0}^{m-1}\omega^{2k}
  + \frac{|w|^2}{2}\cdot m
  + \frac{w^2}{4}\sum_{k=0}^{m-1}\bar\omega^{2k}.
\]
The key fact about roots of unity (the \emph{discrete orthogonality
relation}) is
\begin{equation}
\label{eq:roots-of-unity-sum}
  \sum_{k=0}^{m-1} \omega^{jk} =
  \begin{cases}
    m, & \omega^j = 1 \text{ (i.e.\ } m \mid j), \\
    0, & \text{otherwise},
  \end{cases}
\end{equation}
which follows from the same finite geometric series identity as in
Proof~1: if $\omega^j\ne1$ then
$\sum_k(\omega^j)^k=\frac{(\omega^j)^m-1}{\omega^j-1}=\frac{1-1}{\omega^j-1}=0$,
since $(\omega^j)^m=(\omega^m)^j=1^j=1$. Applying
\eqref{eq:roots-of-unity-sum} with $j=2$: if $m\ge 3$, then $\omega^2\ne
1$ (as $2$ is not a multiple of $m$ for $m\ge3$), so $\sum_k\omega^{2k}=0$
and likewise $\sum_k\bar\omega^{2k}=0$. The two outer sums vanish, leaving
only the middle term:
\[
  \sum_{k=0}^{m-1}|\inner{w}{f_k}|^2 = \frac{|w|^2}{2}\cdot m = \frac{m}{2}|w|^2.
\]
Since $|w|^2 = \norm{w}^2$ under the identification $\C\cong\R^2$, this is
exactly the tight frame identity with $A=m/2$, for every $w\in\R^2$.
\end{proof}

\begin{remark*}[The case $m=2$]
\label{rmk:m-equals-2}
When $m=2$, $\omega^2=e^{4\pi i/2}=e^{2\pi i}=1$, so the cancellation
argument above does not apply directly to the $j=2$ case of
\eqref{eq:roots-of-unity-sum}---indeed $\sum_{k=0}^1\omega^{2k}=2\ne 0$ in
that case. Let us check $m=2$ by hand instead: $f_0=(1,0)$,
$f_1=(-1,0)$ (antipodal points, not really a ``polygon'' in the usual
sense). Direct computation gives $S = f_0f_0^T+f_1f_1^T =
\begin{pmatrix}1&0\\0&0\end{pmatrix}+\begin{pmatrix}1&0\\0&0\end{pmatrix}
=\begin{pmatrix}2&0\\0&0\end{pmatrix}$, which is \emph{not} a multiple of
the identity ($\lambda_{\min}=0\ne\lambda_{\max}=2$). So $m=2$ is
genuinely an exception: two antipodal points do not form a tight frame
(indeed they do not even span $\R^2$!). The theorem's hypothesis $m\ge2$
should really be sharpened to $m\ge3$ for the conclusion to hold without
further case-checking; we keep $m \ge 3$ as the operative hypothesis from
here on, consistent with every genuine ``polygon'' (which needs at least
three vertices).
\end{remark*}

\begin{keyidea}
\textbf{Every regular $m$-gon, $m\ge3$, is a tight frame for $\R^2$},
with frame bound $A = m/2$. The proof reduces to the discrete
orthogonality relation for roots of unity,
$\sum_{k=0}^{m-1}\omega^{jk}=0$ unless $m\mid j$. This is the rigorous,
general version of the cross-term cancellation you first saw by hand in
Example~2.2.5 for $m=3$, and it explains why the cancellation was never a
coincidence: the angles $2\pi k/m$ are always symmetric enough, for any
$m\ge3$, to force the relevant trigonometric sums to vanish.
\end{keyidea}

\begin{remark*}
The value $A=m/2$ matches the trace-formula prediction from
Chapter~3: each $f_k$ has unit norm, so $\tr(S)=\sum_k\norm{f_k}^2=m$;
since the frame is tight in $\R^2$ ($n=2$), tightness forces
$\tr(S)=An=2A$, giving $A=m/2$ exactly as derived above. This consistency
check is a good habit to develop: whenever a tight unit-norm frame
construction is proposed, the trace formula instantly predicts what $A$
\emph{must} be, before any cancellation argument is carried out.
\end{remark*}

\begin{examplebox}[title={Example 4.1: Checking the formula against Chapter 3}]
\label{ex:check-polygon-formula}
Theorem~\ref{thm:regular-polygon-tight} gives $A=m/2$. For $m=3$ (the
triangular frame, Chapter~3's Example~3.2): $A=3/2$, matching the value
$A=B=\tfrac32$ computed there from the explicit matrix $S=\tfrac32I$.
For $m=4$ (the square vertices, Chapter~3's Example~3.6): $A=4/2=2$,
matching the computed value $A=B=2$. Both match exactly, as they must.
\end{examplebox}

\subsection{Why Roots of Unity? A Wider Lens}
\label{subsec:why-roots-of-unity}

It is worth pausing to explain why the roots-of-unity proof is the
``right'' one, beyond simply being shorter. The vectors $f_k = \omega^k$
(under $\C\cong\R^2$) are the orbit of the single vector $f_0=1$ under
repeated rotation by the angle $2\pi/m$. Any construction with this kind
of \emph{cyclic symmetry}---an orbit under a rotation (or more generally,
a unitary matrix) of finite order $m$---will produce cancellation
identities of exactly this type, because the discrete orthogonality
relation~\eqref{eq:roots-of-unity-sum} depends only on the algebraic fact
that $\omega$ is an $m$-th root of unity, not on any special geometric
feature of $\R^2$. This is the first hint of a theme that recurs
throughout frame theory: \textbf{highly symmetric frames are often tight
frames, and group-theoretic symmetry is a systematic tool for
constructing and verifying tightness}, far more powerful than checking
cancellation by hand for each new example. We will see a more general
version of this principle (orbits of finite groups acting on $\R^n$ or
$\C^n$) in Chapter 10 when we discuss equiangular and Grassmannian frames.

\section{Constructing Parseval Frames by Projection}
\label{sec:parseval-from-projection}

We now turn to the second promise of this chapter: explicit examples of
Parseval frames ($A=B=1$) that are not orthonormal bases. Chapter~2
observed this must be possible in principle (a Parseval frame need not be
linearly independent) but gave no construction. Here is a general one.

\begin{theorem}[Projection of an orthonormal basis is Parseval]
\label{thm:projection-parseval}
Let $\{e_1,\ldots,e_N\}$ be an orthonormal basis of $\F^N$, let
$V\subset\F^N$ be an $n$-dimensional subspace ($n \le N$), and let
$P:\F^N\to V$ be the orthogonal projection onto $V$. Then
$\{Pe_1,\ldots,Pe_N\}$ is a Parseval frame for $V$ (viewed as an
$n$-dimensional inner product space in its own right).
\end{theorem}

\begin{proof}
For any $x \in V$, since $P$ is the orthogonal projection onto $V \ni x$,
we have $Px = x$ and $P^*=P=P^2$ (standard properties of orthogonal
projections: self-adjoint and idempotent). Using
$\inner{x}{Pe_k}_{V} = \inner{Px}{Pe_k} = \inner{P^*Px}{e_k} =
\inner{Px}{e_k} = \inner{x}{e_k}$ (the inner product on $V$ is just the
restriction of the inner product on $\F^N$), we compute
\[
  \sum_{k=1}^N |\inner{x}{Pe_k}|^2
  = \sum_{k=1}^N |\inner{x}{e_k}|^2
  = \norm{x}^2,
\]
where the last equality is Parseval's identity (Theorem~1.2.9) for the
orthonormal basis $\{e_k\}$ of $\F^N$, applied to $x \in V \subset \F^N$.
This shows $\{Pe_k\}_{k=1}^N$ satisfies the Parseval identity for every
$x \in V$, i.e.\ it is a Parseval frame for $V$.
\end{proof}

\begin{remark*}
This construction is extremely flexible: start with \emph{any}
orthonormal basis of \emph{any} larger space $\F^N$, project down onto
\emph{any} subspace $V$, and the projected vectors automatically form a
Parseval frame for $V$---regardless of how the subspace sits inside
$\F^N$. Since $N$ can be made as large as we like relative to
$n=\dim V$, this produces Parseval frames with arbitrarily high
redundancy $N/n$. Crucially, the projected vectors $\{Pe_k\}$ are
generally \emph{not} linearly independent once $N>n$ (there are $N>n$ of
them in an $n$-dimensional space), and they are typically not even
distinct or nonzero, so this is genuinely different from an orthonormal
basis while retaining the exact Parseval identity.
\end{remark*}

\begin{examplebox}[title={Example 4.2: A Parseval frame from projecting $\R^3$ onto $\R^2$}]
\label{ex:parseval-projection-example}
Let $N=3$, $\{e_1,e_2,e_3\}$ the standard basis of $\R^3$, and let
$V = \setdef{(x_1,x_2,x_3)\in\R^3}{x_1+x_2+x_3=0}$, the plane through the
origin perpendicular to $(1,1,1)$. This is a $2$-dimensional subspace.
The orthogonal projection onto $V$ is $P = I - \tfrac13 uu^T$ where
$u=(1,1,1)^T/\sqrt3$ is the unit normal to $V$. Computing:
\[
  Pe_1 = e_1 - \tfrac13(1,1,1)^T = \tfrac13(2,-1,-1)^T, \quad\text{and similarly for } Pe_2, Pe_3
\]
by symmetry. The three vectors $Pe_1, Pe_2, Pe_3$ all have the same
length $\sqrt{(2/3)^2+(1/3)^2+(1/3)^2}=\sqrt{6}/3$ and, viewed inside the
$2$-dimensional plane $V$, turn out to be at $120^\circ$ to one another:
this is, geometrically, exactly the triangular frame of
Chapter~1 (up to rotation and an overall scale factor), now arrived at
by an entirely different route---projection from $\R^3$, rather than
direct placement in $\R^2$. By Theorem~\ref{thm:projection-parseval},
$\{Pe_1,Pe_2,Pe_3\}$ is automatically Parseval ($A=B=1$) for $V$, even
though the three vectors are linearly dependent (they live in a
$2$-dimensional space) and are not the standard basis of anything. This
example is verified numerically in Lab Exercise~4.2 below.
\end{examplebox}

\section{A Full Geometric Characterization of Tightness in $\R^2$}
\label{sec:geometric-characterization}

The regular polygon theorem shows that rotational symmetry is
\emph{sufficient} for tightness. It is natural to ask whether it is
\emph{necessary}---must a tight frame in $\R^2$ look like a regular
polygon? The answer is no, and understanding why illuminates the true
geometric content of tightness.

\begin{proposition}[Tight frames in $\R^2$ need not be regular polygons]
\label{prop:tight-not-regular}
There exist tight frames for $\R^2$ whose vectors do not have equal
norms and are not symmetrically arranged.
\end{proposition}

\begin{proof}
We exhibit one. Let $g_1=(1,0)$, $g_2=(0,1)$, $g_3=(1,0)$, $g_4=(0,1)$
(each standard basis vector repeated twice). Then
\[
  S = \sum_{i=1}^4 g_ig_i^T = 2e_1e_1^T + 2e_2e_2^T = 2I,
\]
so this is tight with $A=2$, yet the four vectors are not at equal
angles to one another (two coincide with $e_1$, two with $e_2$; the
angle between the two \emph{distinct} directions present is $90^\circ$,
not the $120^\circ$ a genuine $4$-fold symmetric picture would need) ---
this tight frame is not the vertex set of any regular polygon, and is
simply two copies of an orthonormal basis stacked together.

For a less degenerate example, with vectors of \emph{unequal} norm, fix
any $a, b, c, d > 0$ satisfying $a^2+b^2 = c^2+d^2$ (for instance
$a=1.5,\,b=0.5,\,c=1,\,d=\sqrt{1.5^2+0.5^2-1^2}=\sqrt{1.5}$), and set
\[
  g_1 = (a, 0), \quad g_2 = (-b, 0), \quad g_3 = (0, c), \quad g_4 = (0, -d).
\]
Since $g_1, g_2$ lie on the $x$-axis and $g_3, g_4$ lie on the $y$-axis,
the frame operator is automatically diagonal:
\[
  S = (a^2+b^2)\,e_1e_1^T + (c^2+d^2)\,e_2e_2^T,
\]
which equals $A\cdot I$ for $A = a^2+b^2 = c^2+d^2$ by our choice of
constants. Here the four vectors visibly have four different lengths
($a=1.5,\,b=0.5,\,c=1,\,d\approx1.225$ in the example above) and are
certainly not the vertices of any regular polygon, yet the frame is
exactly tight. This proves the proposition. (A fully asymmetric example
--- vectors at generic, non-axis-aligned angles, with unequal norms and
no special structure at all --- is also possible, but is more easily
found numerically than by hand; Lab Exercise~4.4 asks you to search for
one.)
\end{proof}

\begin{remark*}
The correct general statement (which we will prove rigorously in
Chapter~9 using the \emph{frame potential}) is this: a set of $m$
unit-norm vectors in $\R^n$ is a tight frame if and only if it is a
\emph{critical point}---in fact a global minimizer---of the total
pairwise-correlation energy $\sum_{i\ne j}|\inner{f_i}{f_j}|^2$. Regular
polygons are tight because their symmetry forces them to be minimizers,
but they are far from the only minimizers once $m$ or $n$ grow: there is
generally a large family of inequivalent tight configurations. Tightness
is therefore better understood as an \emph{energy-minimization}
(equivalently, \emph{balance}) condition than as a symmetry condition;
symmetry is a convenient sufficient mechanism for achieving it; but it is
not the only mechanism, and Chapter~9 develops the energy perspective in
full.
\end{remark*}

\section{Computing and Verifying Tight Frames in NumPy}
\label{sec:numpy-tight-frames}

We now translate the regular polygon theorem and the projection
construction into code, building tools to construct and verify tight
frames numerically.

\begin{lstlisting}[caption={Regular polygon frames and the projection construction}]
import numpy as np

def regular_polygon_frame(m):
    """
    Return the n=2 synthesis matrix (shape (2, m)) for the
    vertices of a regular m-gon inscribed in the unit circle.
    """
    angles = 2 * np.pi * np.arange(m) / m
    return np.vstack([np.cos(angles), np.sin(angles)])

def frame_operator(F):
    return F @ F.T

def exact_frame_bounds(F):
    S = frame_operator(F)
    eigenvalues, eigenvectors = np.linalg.eigh(S)
    return float(eigenvalues[0]), float(eigenvalues[-1]), eigenvalues, eigenvectors

def is_tight(F, tol=1e-9):
    A, B, _, _ = exact_frame_bounds(F)
    return (B - A) < tol * max(1.0, B)

# ---------------------------------------------------------------
# Theorem 4.3.1: verify A = m/2 for many regular polygons
# ---------------------------------------------------------------
print("=== Regular polygon frames ===")
for m in range(3, 13):
    F = regular_polygon_frame(m)
    A, B, eigs, _ = exact_frame_bounds(F)
    predicted_A = m / 2
    print(f"m={m:2d}:  A={A:.6f}  B={B:.6f}  "
          f"predicted m/2={predicted_A:.4f}  "
          f"tight={is_tight(F)}  match={np.isclose(A, predicted_A)}")

# ---------------------------------------------------------------
# Theorem 4.4.1: Parseval frame via projection (Example 4.2)
# ---------------------------------------------------------------
def projection_parseval_frame(N, V_basis):
    """
    Given an orthonormal basis (columns) of a subspace V of R^N,
    return the Parseval frame for V obtained by projecting the
    standard basis of R^N onto V, expressed in V's own coordinates.

    V_basis: ndarray of shape (N, n), orthonormal columns spanning V.
    Returns: ndarray of shape (n, N) -- N frame vectors for the
             n-dimensional space V, in V's coordinate system.
    """
    # P = V_basis @ V_basis.T projects R^N onto V (as vectors in R^N).
    # Expressed in V's own n-dimensional coordinates, the projection
    # of e_k is simply the k-th row of V_basis (i.e. V_basis.T @ e_k).
    return V_basis.T  # shape (n, N): columns are the N frame vectors

# Example 4.2: V = {x : x1+x2+x3=0} subset of R^3
u = np.array([1., 1., 1.]) / np.sqrt(3)   # unit normal to V
# Build an orthonormal basis of V (2-dimensional):
v1 = np.array([1., -1., 0.]) / np.sqrt(2)
v2 = np.array([1., 1., -2.]) / np.sqrt(6)
V_basis = np.column_stack([v1, v2])        # shape (3, 2)

F_proj = projection_parseval_frame(3, V_basis)   # shape (2, 3)
A_proj, B_proj, eigs_proj, _ = exact_frame_bounds(F_proj)
print("\n=== Parseval frame via projection (Example 4.2) ===")
print("Frame vectors (in V's 2D coordinates):\n", F_proj.round(4))
print(f"A = {A_proj:.6f}, B = {B_proj:.6f}  (expect A=B=1)")
\end{lstlisting}

Running the first loop confirms Theorem~\ref{thm:regular-polygon-tight}
for $m=3,\ldots,12$ in a single pass: every regular polygon is tight,
with $A$ matching $m/2$ to numerical precision. The second block
constructs the projection-based Parseval frame of Example~4.2 from
scratch and confirms $A=B=1$ exactly.

\begin{remark*}
A subtlety in the \texttt{projection\_parseval\_frame} function: rather
than literally constructing the $N\times N$ projection matrix $P$ and
then re-expressing each $Pe_k\in\R^N$ inside the lower-dimensional
subspace $V$, we exploit a shortcut. If $V\_basis$ has orthonormal
columns $v_1,\ldots,v_n$ spanning $V$, then the coordinates of $Pe_k$
\emph{within} $V$ (relative to the basis $\{v_j\}$) are exactly
$(\inner{e_k}{v_1},\ldots,\inner{e_k}{v_n}) = $ the $k$-th row of
\texttt{V\_basis}. This avoids ever forming the $N\times N$ matrix $P$
explicitly, which matters for efficiency when $N$ is large, and is worth
internalizing as a general pattern: orthogonal projection coordinates are
just inner products against an orthonormal basis of the target subspace.
\end{remark*}

\section{Chapter Summary}
\label{sec:summary-ch4}

\begin{itemize}
  \item Tightness is equivalent to a \textbf{resolution of the
        identity}: $\sum_i f_if_i^* = AI$
        (Proposition~\ref{prop:tightness-resolution}), which entrywise
        says diagonal sums equal $A$ and off-diagonal sums vanish.
  \item \textbf{Theorem~\ref{thm:regular-polygon-tight}}: every regular
        $m$-gon ($m\ge3$) inscribed in the unit circle is a tight frame
        for $\R^2$ with bound $A=m/2$, matching the trace-formula
        prediction $A=\tr(S)/n=m/2$. The proof reduces to the
        \textbf{discrete orthogonality relation for roots of unity},
        $\sum_{k=0}^{m-1}\omega^{jk}=0$ unless $m\mid j$---the rigorous,
        general mechanism behind the $120^\circ$ cross-term cancellation
        first seen by hand in Example~2.2.5.
  \item The case $m=2$ is a genuine exception (antipodal points do not
        even span $\R^2$); the theorem requires $m\ge3$.
  \item \textbf{Theorem~\ref{thm:projection-parseval}}: projecting any
        orthonormal basis of a larger space $\F^N$ onto a subspace $V$
        always produces a Parseval frame for $V$, regardless of how $V$
        sits inside $\F^N$. This gives explicit, flexible examples of
        Parseval frames that are not orthonormal bases, fulfilling the
        promise of Chapter~2.
  \item Tightness is best understood as an \textbf{energy-balance}
        condition (to be made precise via the frame potential in
        Chapter~9), not merely a symmetry condition: symmetric
        configurations (regular polygons, roots of unity) are a
        convenient sufficient mechanism, but not the only tight
        configurations.
\end{itemize}

\section{Lab Exercises}
\label{sec:lab-ch4}

\begin{exercisebox}[title={Lab Exercise 4.1: The Full Regular Polygon Table}]
Implement \texttt{regular\_polygon\_frame(m)} and
\texttt{exact\_frame\_bounds(F)} from
Section~\ref{sec:numpy-tight-frames}. Build a table of $A$, $B$, and
$\tr(S)$ for $m=3,4,\ldots,20$. Verify in every row that (i) $A=B$
(tight), (ii) $A=m/2$, and (iii) $\tr(S)=m$. Then check $m=2$ separately
and confirm it breaks the pattern, consistent with
Remark~\ref{rmk:m-equals-2}.
\end{exercisebox}

\begin{exercisebox}[title={Lab Exercise 4.2: Verifying the Projection Construction}]
Implement \texttt{projection\_parseval\_frame} from
Section~\ref{sec:numpy-tight-frames} and reproduce Example~4.2 exactly.
Then construct a second example: let $V\subset\R^4$ be the
$2$-dimensional subspace spanned by an orthonormal basis of your choosing
(e.g.\ via \texttt{np.linalg.qr} applied to two random vectors in
$\R^4$), and verify that projecting the standard basis $\{e_1,e_2,e_3,e_4\}$
of $\R^4$ onto $V$ produces a $4$-vector Parseval frame for the
$2$-dimensional space $V$. Confirm $A=B=1$ numerically.
\end{exercisebox}

\begin{exercisebox}[title={Lab Exercise 4.3: Partial Regular Polygons Are Not Tight}]
Theorem~\ref{thm:regular-polygon-tight} requires using \emph{all} $m$
vertices of the regular polygon. Investigate what happens with a
\emph{partial} set: for $m=8$, compute the frame operator and exact
bounds for the first $m'=2,3,4,5,6,7$ vertices only (i.e.\
$f_k=(\cos(2\pi k/8),\sin(2\pi k/8))$ for $k=0,\ldots,m'-1$, omitting the
rest). For which values of $m'$ (if any) is the resulting frame still
tight? Reconcile your numerical finding with the roots-of-unity proof:
which step of Proof~2 specifically breaks down when the sum is only over
a partial set of $k$ values?
\end{exercisebox}

\begin{exercisebox}[title={Lab Exercise 4.4: Hunting for Asymmetric Tight Frames}]
Proposition~\ref{prop:tight-not-regular} claims tight frames need not be
regular polygons. Investigate numerically: starting from the repeated
basis vectors example in the proof
($g_1=g_3=(1,0)$, $g_2=g_4=(0,1)$), perturb each vector slightly by a
small random rotation (a few degrees) and a small random rescaling
(within $\pm10\%$ of unit length), and check whether the perturbed frame
is still (approximately) tight. Report the ratio $B/A$ before and after
perturbation. Then try to manually construct (by trial and error, or
using \texttt{scipy.optimize}) a 4-vector frame in $\R^2$ that is exactly
tight but visibly \emph{not} a regular polygon and \emph{not} two
repeated orthonormal bases --- i.e.\ four genuinely distinct directions,
no two equal or antipodal, with no obvious rotational symmetry, that
still satisfies $S=AI$. (Hint: you have 4 vectors $\times$ 2 coordinates
$=8$ real degrees of freedom, and the tight-frame condition $S=AI$ in
$\R^2$ imposes only 3 real constraints---2 diagonal entries equal, 1
off-diagonal entry zero--- so a $5$-real-dimensional family of solutions
should exist; you only need to find one element of it.)
\end{exercisebox}

\begin{exercisebox}[title={Lab Exercise 4.5 (Optional, Challenge): Roots of Unity in $\C^n$}]
The roots-of-unity proof technique generalizes beyond $\R^2$. Consider
the $n$-dimensional complex vectors
\[
  f_k = \frac{1}{\sqrt{n}}\bigl(1, \omega^k, \omega^{2k}, \ldots, \omega^{(n-1)k}\bigr) \in \C^n,
  \qquad k = 0, 1, \ldots, n-1,
\]
where $\omega = e^{2\pi i/n}$ (these are the columns of the
\emph{discrete Fourier transform matrix}, up to normalization).
\begin{enumerate}[label=(\alph*)]
  \item For $n=4$, construct these vectors numerically (using complex
        NumPy arrays, \texttt{dtype=complex}) and verify that
        $\{f_k\}_{k=0}^{n-1}$ is an orthonormal basis of $\C^n$ (hence,
        trivially, a tight frame with $A=B=1$) by checking
        $\inner{f_j}{f_k}=\delta_{jk}$ for all $j,k$ (recall: for
        complex vectors, use \texttt{np.vdot} as discussed in
        Chapter~1's remark on conventions).
  \item Now take only the first $n-1$ of these vectors (drop $f_{n-1}$)
        for $n=4$. Is the resulting $3$-vector set still a tight frame
        for $\C^4$? (It cannot be, by Theorem~2.2.1, for a reason you
        should be able to state immediately --- what is it?)
  \item This construction --- equally spaced points on the unit circle
        in each coordinate, built from roots of unity --- is the
        starting point for an entire family of frames called
        \emph{harmonic frames}, which we will study properly in
        Chapter~10 alongside equiangular frames. For now, simply record
        your numerical observations as a preview.
\end{enumerate}
\end{exercisebox}

\newpage
\section{Supplementary Exercises}
\label{sec:extra-ch4}

\subsection*{Theoretical Exercises}
\begin{theoexercises}
	\item Prove that if $\{f_i\}_{i=1}^m$ is a tight frame for $\R^n$ with
	bound $A$, then taking $k$ copies of each vector (each vector
	repeated $k$ times, for a total of $km$ vectors) is also tight,
	with bound $kA$.
	
	\item Prove that if $\{f_i\}_{i=1}^m$ is a tight frame for $\R^n$ with
	bound $A_1$ and $\{g_j\}_{j=1}^{m'}$ is a tight frame for $\R^n$
	with bound $A_2$ (independently, on a disjoint index set), then
	the combined set $\{f_i\}\cup\{g_j\}$ is a tight frame for $\R^n$
	with bound $A_1+A_2$ --- \emph{regardless of whether $A_1=A_2$}.
	
	\item Prove directly, via a symmetric summation argument (not the
	roots-of-unity technique of Section~4.3), that
	$\{\pm e_1,\pm e_2,\pm e_3\}$ is a tight frame for $\R^3$, and
	determine its bound $A$.
	
	\item Prove that adding a single unit vector $v$ to an orthonormal
	basis $\{e_1,\ldots,e_n\}$ of $\R^n$ ($n\ge2$) never produces a
	tight frame, for any choice of $v$. (Hint: the resulting frame
	operator is $I+vv^T$; compare its rank structure to a scalar
	multiple of $I$.)
	
	\item Generalize the axis-aligned construction of
	Proposition~4.5.1 to $\R^3$: using six vectors, two along each
	coordinate axis with possibly unequal norms $a,b$ (on the first
	axis), $c,d$ (second), $e,f$ (third), state and prove the precise
	condition on these six numbers for the configuration to be tight.
	
	\item Using Theorem~8.3.1 (frame potential $\ge m^2/n$, with equality
	iff tight), give an alternative proof that the regular $m$-gon is
	a tight frame: compute its frame potential directly using a
	double-angle identity, compare to the lower bound $m^2/n$ with
	$n=2$, and conclude.
\end{theoexercises}

\subsection*{Computational Exercises}
\begin{compexercises}
	\item Numerically verify T1 for the triangular frame with $k=2,3,4$
	copies, confirming the bound scales as $kA$.
	
	\item Numerically verify T2 using the triangular frame ($A=1.5$) and
	the square ($A=2$) as the two tight frames being combined, and
	confirm the union is tight with bound $3.5$.
	
	\item Compute the frame operator of $\{\pm e_1,\pm e_2,\pm e_3\}$
	directly and confirm the eigenvalues found in T3.
	
	\item For $n=2$ and $n=3$, add a random unit vector to the standard
	orthonormal basis and confirm numerically that the result is
	never tight (T4), across at least $10$ random choices of the
	extra vector in each dimension.
	
	\item Construct the $6$-vector $\R^3$ configuration from T5 with a
	specific numeric target bound (e.g.\ $A=5$), verify tightness
	directly, and confirm the frame operator is exactly $5I$.
\end{compexercises}
